% chktex-file 45
\section{Preliminaries}\label{sec: preliminaries}
%%%%%%%%%%%%%%%%%%%%%%%%%%%%%%

We first introduce the necessary notation.
The distance between points $x$ and $y$ is denoted by $|xy|$, $d(x,y)$ or $\rho(x,y)$.
The open ball with center $x$ and radius $r$ is denoted by $U_r(x)$ or $U(x,r)$, and the closed one by $B_r(x)$ or $B(x,r)$.

\begin{define}
  Let $A,B$ be non-empty subsets of a metric space $X$.
The \emph{Hausdorff distance} is the quantity
  \begin{equation*}
      \di{H}{X}(A, B) = \inf \Bigl\{r\: A \subset U_r(B)\, \&\, B \subset U_r (A)\Bigr\}.
  \end{equation*}
\end{define}

\begin{define}
  Let $A,B,X$ be metric spaces.
If $A$ is isometric to $\vA$ and $B$ is isometric to $\vB$, where $\vA$ and $\vB$ are subspaces of $X$, then the triple $(\vA, \vB, X)$ is called a \emph{realization of the pair} $(A, B)$.
\end{define}
      

\begin{define}\label{def: gh-realization}
  The \emph{Gromov--Hausdorff distance} between two metric spaces $A$, $B$ is the infimum of the Hausdorff distances over all realizations of the pair $(A,B)$.
In other words,
  \begin{equation*}
      \dGH{}(A,B) = \inf\bigl\{r: \text{there exists a realization } (\vA,\vB,X) \text{ of the pair $(A,B)$ with }\di{H}{X} (\vA,\vB) \leq r \bigr\}.
  \end{equation*}
\end{define}
As stated in the introduction, the class $\GH$ consists of representatives: one space with a fixed metric is chosen from each isometry class of metric spaces. The notation $X \in \GH$ means such a particular metric space, not an isometry class.
We now give another, simpler definition of the Gromov--Hausdorff distance.
\begin{define}
  Let $X$, $Y$ be metric spaces and let $\sigma \subseteq X \x Y$.
Its \emph{distortion} is defined by
  $$
  \dis \sigma = \sup \Bigl\{\bigl|\dist{X}(a,a')-\dist{Y}(b,b')\bigr|\:\, (a, b) \text{ and } (a', b') \in \sigma \Bigr\}.
  $$
  It is easy to see that if $\sigma'\subseteq\sigma$, then $\dis\sigma'\le\dis\sigma$, since the $\sup$ over a subset does not exceed the $\sup$ over the containing set.
\end{define}
The distortion of a map is the distortion of its functional correspondence:
$$
\dis f = \sup_{a, a' \in X}\bigl|\dist{X}(a,a') - \dist{Y}(f(a), f(a'))\bigr|
$$
\begin{define}
  A \emph{correspondence} between two sets $A$ and $B$ is a subset $R \subset A \x B$ such that for any $a \in A$ and $b \in B$ there exist $\va \in A$ and $\vb \in B$ with $(a,\vb)$, $(\va, b)$ in $R$.
\end{define}
We write $aRb$ to say that $a$ and $b$ are related by $R$, and denote the set of all correspondences between metric spaces $A$, $B$ by $\RR(A,B)$.
For $a \in A$ we write $R(a) = \{b \in B : (a, b) \in R\}$; similarly $R^{-1}(b) = \{a \in A : (a, b) \in R\}$ for $b \in B$.

\begin{define}\label{def: simplex}
  A metric space with $n$ points in which all non-zero distances are equal is called the \emph{$n$-point simplex} $\Delta_n$. If these distances are equal to $\lambda > 0$, we denote it by $\lambda\Delta_n$.
\end{define}


  
\begin{proposition}[\cite{BurBurIva}]\label{theorem:main_formula}
  For any metric spaces $A$ and $B$,
  \begin{equation*}
      2\dGH{}(A,B) = \underset{R \in \RR(A,B)}{\inf} \dis R.
  \end{equation*}
\end{proposition}


\begin{example}
  The simplest example of a correspondence $R \in \RR(X,Y)$ comes from a pair of maps $f \: X \to Y$ and $g \: Y \to X$, viewed as functional correspondences, by taking their ``union'': $R_{f,g} = f \cup g^{-1}$.
\end{example}
\begin{define}
  The \emph{codistortion} $\codis(f,g)$ is the quantity
  $$
  \codis(f,g)=\sup_{x\in X,\,y\in Y}\Bigl|\,\bigl|x\,g(y)\bigr|-\bigl|f(x)\,y\bigr|\,\Bigr|.
  $$
\end{define}
For the correspondence $R_{f,g}$ from the previous example,
$$
\dis R_{f,g} = \max\bigl\{\dis f,\,\dis g,\,\codis(f,g)\bigr\}:
$$
pairs of points in $R_{f,g}$ are of three kinds --- both taken from $f$, both from $g^{-1}$, or one from each --- and the suprema over them give exactly $\dis f$, $\dis g$ and $\codis(f,g)$.

\begin{proposition}[\cite{TuzCont}]~\label{prop: gh via maps}
  For any $X,Y\in\GH$,
  $$
  \dGH{}(X,Y) =  
   \frac12\,\inf_{\substack{f\:X\to Y\\ g\:Y\to X}}\, \dis R_{f,g} = 
  \frac12\,\inf_{\substack{f\:X\to Y\\ g\:Y\to X}}\,\sup\bigl\{\dis f,\,\dis g,\,\codis(f,g)\bigr\}.
  $$
\end{proposition}

The formula shows that the larger the class of morphisms, the smaller the infimum: if the morphisms of a subcategory $\mathcal{M}_1$ lie among the morphisms of $\mathcal{M}_2$, then $\dGH{\FF_2}\le\dGH{\FF_1}$ (formally: Proposition~\ref{prop:FFmonotone}).

\begin{proposition}[\cite{TuzCont}]~\label{prop: triangle inequality}
Let $X$, $Y$ and $Z$ be three metric spaces, and let $f: X \to Y$, $g: Y \to X$, $h: Y \to Z$, $k: Z \to Y$ be maps, denoted compactly as $X\torllim^f_gY\torllim^h_kZ$ (the upper indices are the ``forward'' maps, the lower ones the ``backward'' maps). Then $\codis(h\c f,g\c k)\le\codis(f,g)+\codis(h,k)$.
\end{proposition}
We introduce the notation $\mathbb{R}_+^n = \{(x_1, \ldots, x_n) \in \mathbb{R}^n : x_1 \geq 0\}$.

\begin{define}
    Let $M$ be a Hausdorff space with a countable base.
    A family of pairs $\{(U_i,\varphi_i)\}$, where the $U_i$ are open sets
    covering $M$ and each $\varphi_i$ is a homeomorphism of $U_i$ onto an open
    subset of $\mathbb{R}^n_+$, is called a $C^k$-atlas if
    the maps
    $$
    \varphi_i \circ \varphi_j^{-1} \bigl|_{\varphi_j(U_i \cap U_j)}
    $$
    are of class $C^k$ for all $i,j$.

    A $C^k$-structure on $M$ is a $C^k$-atlas maximal by inclusion.
    The pair $(M, C^k\textup{-structure})$ is called a $C^k$-manifold
    with boundary (possibly empty).
\end{define}
\begin{comm}
  Smoothness at the boundary is understood as follows. A function defined on a subset of the half-space $\R^n_+$ that is open in $\R^n_+$ is called $C^k$-smooth if it extends to a $C^k$-smooth function on some open subset of $\R^n$. This is consistent with the transition maps of the atlas and defines the class $C^k$ at boundary points.
\end{comm}
\begin{comm}\label{comm: paracompact}
A $C^k$-manifold (with or without boundary) is paracompact.
\end{comm}
\begin{comm}
  In this work, by a ``manifold'' we mean both ordinary manifolds and manifolds with boundary.
\end{comm}
\begin{define}
    Let $M$ and $N$ be $C^k$-smooth manifolds and let
    $f : M \to N$ be a map.

    The map $f$ is called \emph{$C^r$-smooth} ($0 \le r \le k$) if
    for each point $x \in M$ there are a chart $(U, \varphi)$ of $M$
    with $x \in U$ and a chart $(V, \psi)$ of $N$ with $f(x) \in V$, such that
    $f(U) \subset V$ and the map
    $$
    \psi \circ f \circ \varphi^{-1} :
    \varphi(U) \to \psi(V)
    $$
    is a $C^r$-smooth map between open subsets of
    $\mathbb{R}^n$ and $\mathbb{R}^m$.
\end{define}

In the classical literature on Gromov triples one usually restricts to probability or finite measures, but we work with a larger class of measures.
The reason is the following: submanifolds of $\R^n$ are often unbounded, so the measure induced on them by the Lebesgue measure is not finite (but it is boundedly finite).
\begin{define}
  An ordered triple $(X, d, \mu)$, where
  \begin{enumerate}
    \item $(X,d)$ is a metric space,
    \item $\mu$ is a Borel boundedly finite measure (all balls, or equivalently all bounded sets, have finite measure),
  \end{enumerate}
  is called a \emph{Gromov metric triple}.
  Recall that a map $F \: X \to Y$ between spaces with Borel sigma-algebras is called \emph{measurable} if the preimage $F^{-1}(B)$ of any Borel set $B \subset Y$ is Borel in $X$; for a measurable $F$, the pushforward of the measure is defined by $(F_* \mu_X)(B) = \mu_X(F^{-1}(B))$ for all Borel sets $B \subset Y$.
  Let $X = (X, \dist{X}, \mu_X)$ and $Y = (Y, \dist{Y}, \mu_Y)$ be Gromov metric triples. A map $F \: X \to Y$ is called an \emph{isomorphism} of Gromov metric triples if $F$ is an isometry of the underlying metric spaces that preserves the measure: $F_* \mu_X = \mu_Y$. An isometry is continuous and hence Borel measurable, so $F_* \mu_X$ is well defined.
  We denote by $\GH_{mm}$ the \emph{class of Gromov metric triples}, choosing one representative in each isomorphism class.
  In this work, the distance between Gromov metric triples is computed as the distance between their underlying metric spaces.
\end{define}
